\documentclass[10pt,a4paper]{amsart}
\usepackage[margin=19mm]{geometry}
\usepackage{amsmath,amssymb,mathtools}
\usepackage[scr=rsfs,cal=euler]{mathalfa}
\usepackage{xfrac}
\usepackage[unicode,hidelinks,bookmarks=false]{hyperref}
\newcommand{\ie}{i.e.\ }
\newcommand{\cf}{cf.\ }
\newcommand{\resp}{respectively\ }
\newcommand{\Subsection}{\subsection}
\providecommand{\nobreakdash}{\nobreak\hbox{-}}
\setlength{\emergencystretch}{2em}
\multlinegap0pt
\newtheorem{lemma}{Lemma}
\title{A discrete mean-value theorem for the higher derivatives of the Riemann zeta function}
\author{Christopher Hughes, Andrew Pearce-Crump}
\date{Source: arXiv:2106.03005v3 (CC BY 4.0)}
\begin{document}
\maketitle

\noindent\emph{Source and licence.} A discrete mean-value theorem for the higher derivatives of the Riemann zeta function. Version arXiv:2106.03005v3 (CC BY 4.0), distributed by arXiv under the Creative Commons Attribution 4.0 International licence, http://creativecommons.org/licenses/by/4.0/, as recorded in the arXiv licence metadata for this exact version and shown on the abstract page https://arxiv.org/abs/2106.03005v3. Full licence notice: \texttt{licenses/arxiv-2106.03005-CC-BY-4.0.txt}.\par\smallskip

\noindent\textit{Editorial scope.} Counted unit: the article's lemma 11 (source0.tex lines 414--420). The article's complete original proof of it is delivered with it (source0.tex lines 422--454, one \texttt{proof} environment of 33 lines). No mathematics is written, repaired or re-derived: the statement and the proof are the article's own lines, in the article's own notation, and no retained line is substituted or re-flowed.  Retained uncounted as the source-local context the proof needs: lines 396--399; lines 405--410.  Nothing was omitted from any retained span, so the package carries no omission record.  No retained line contains a citation, so the wrapper carries no bibliography and no bibliography entry.  No retained line contains a figure environment, an image-inclusion command or drawing code, so no image is delivered and the package declares no asset dependency.  Editorial formatting only, in addition: an AMS \texttt{amsart} preamble that restates the article's own declarations and macros used by the retained lines, namely the environments \texttt{lemma} on separate counters, together with the standard packages those lines need.  The retained environments print in this document as Lemma 1 in order of appearance, whereas the article prints the same items with the numbering of its own full document.  The complete native source of the article as distributed in the arXiv e-print for this exact version is bundled unchanged as \texttt{sources/112322/source0.tex}.
\begin{align}
I &= \frac{(-1)^{n+1}}{2 \pi} \int_{1}^{T} \chi (1 - c - it) \left[ \sum_{k=0}^n \binom{n}{k} \log^{n-k+1} \left( \frac{t}{2 \pi} \right) \zeta^{(k)}(c +it) \right. \notag\\
&+\left. \sum_{k=0}^n \binom{n}{k} \log^{n-k} \left( \frac{t}{2 \pi} \right) \zeta^{(k)}(c+it) \frac{\zeta '}{\zeta} (c +it) \right] \ dt + O(T^{\frac{1}{2}} \log ^{n+2} T). \label{eq:I split}
\end{align}

\begin{lemma}[Gonek]\label{Gonek 5}
Let $\{b_m\}_{m=1}^{\infty}$ be a sequence of complex numbers such that for any $\varepsilon>0$, $b_m \ll m^\varepsilon$. Let $c>1$ be as before and let $k$ be a non-negative integer. Then for $T$ sufficiently large,
\[
\frac{1}{2 \pi} \int_{1}^{T} \left( \sum_{m=1}^{\infty} b_m m^{-c -it} \right) \chi (1-c-it) \log ^k \left(\frac{t}{2 \pi} \right) \ dt = \sum_{1 \leq m \leq \frac{T}{2 \pi}} b_m \log^k m + O(T^{c-\frac{1}{2}} \log ^k T).
\]
\end{lemma}

\begin{lemma}\label{Lemma 11}
The integral $I$ can be written as
\[
I= (-1)^n \sum_{m r \leq \frac{T}{2 \pi}} \Lambda(r) \log^n r + O(T^{\frac{1}{2}}\log^{n+2} T)
\]
where the sum is taken over all integers $m$ and $r$ such that $mr \leq \frac{T}{2\pi}$.
\end{lemma}

\begin{proof}
Before we apply Lemma~\ref{Gonek 5}, we will split the integral $I$ given in \eqref{eq:I split} substituting the Dirichlet series to simplify the argument slightly. To do this, we write
\begin{align*}
I &= \frac{(-1)^{n+1}}{2 \pi } \int_{1}^{T} \chi (1 - c - it) \sum_{k=0}^n \binom{n}{k} \log^{n-k+1} \left( \frac{t}{2 \pi} \right) \sum_{m=1}^{\infty} \frac{(-1)^k \log^k m}{m^{c + it}} \ dt \\
& \quad +\frac{(-1)^{n+1}}{2 \pi } \int_{1}^{T} \chi (1 - c - it) \sum_{k=0}^n \binom{n}{k} \log^{n-k} \left( \frac{t}{2 \pi} \right) \sum_{m=1}^{\infty} \frac{(-1)^{k+1} \log^k m}{m^{c + it}} \sum_{r=1}^{\infty} \frac{\Lambda (r)}{r^{c+it}} \ dt \\
& \quad + O(T^{\frac{1}{2}} \log^{n+2} T) \\
&= J_1 + J_2 +O(T^{\frac{1}{2}} \log^{n+2} T),
\end{align*}
say.

For $J_1$, we have
\begin{align*}
J_1 &= (-1)^{n+1}\sum_{k=0}^n \binom{n}{k} (-1)^k  \frac{1}{2 \pi } \int_{1}^{T} \chi (1 - c - it) \log^{n-k+1} \left( \frac{t}{2 \pi} \right) \sum_{m=1}^{\infty} \frac{ \log^k m}{m^{c + it}} \ dt \\
&= (-1)^{n+1} \sum_{k=0}^n \binom{n}{k} (-1)^k \left( \sum_{m \leq \frac{T}{2 \pi}} \log^k m \log^{n-k+1} m +O(T^{\frac{1}{2}}\log^{n-k+1} T) \right)\\
&= (-1)^{n+1} \sum_{k=0}^n \binom{n}{k} (-1)^k \sum_{m \leq \frac{T}{2 \pi}} \log^{n+1} m + O(T^{\frac{1}{2}}\log^{n+1} T) \\
&=0 + O(T^{\frac{1}{2}}\log^{n+1} T)
\end{align*}
since by the Binomial Theorem,
\[
\sum_{k=0}^n \binom{n}{k} (-1)^k = (1 + (-1))^n = 0.
\]

For $J_2$, we have
\begin{align*}
J_2 &= (-1)^{n+1} \sum_{k=0}^n \binom{n}{k} (-1)^{k+1} \frac{1}{2 \pi } \int_{1}^{T} \chi (1 - c - it)  \log^{n-k} \left( \frac{t}{2 \pi} \right) \sum_{m=1}^{\infty} \frac{ \log^k m}{m^{\sigma + it}} \sum_{r=1}^{\infty} \frac{\Lambda (r)}{r^{c+it}} \ dt \\
&= (-1)^{n+1} \sum_{k=0}^n \binom{n}{k} (-1)^{k+1} \frac{1}{2 \pi } \int_{1}^{T} \chi (1 - c - it)  \log^{n-k} \left( \frac{t}{2 \pi} \right) \sum_{mr=1}^{\infty} \frac{ \Lambda (r) \log^k m }{(mr)^{c + it}} \ dt \\
&= (-1)^{n+1} \sum_{k=0}^n \binom{n}{k} (-1)^{k+1} \left( \sum_{m r \leq \frac{T}{2 \pi}} \Lambda (r) \log^k (m) \log^{n-k} (m r) + O(T^{\frac{1}{2}}\log^{n-k} T)\right)\\
&= (-1)^{n} \sum_{m r \leq \frac{T}{2 \pi}} \Lambda (r) \sum_{k=0}^n \binom{n}{k} (-\log m)^k \log^{n-k} (m r)  + O(T^{\frac{1}{2}}\log^{n} T) \\
&= (-1)^{n} \sum_{m r \leq \frac{T}{2 \pi}} \Lambda (r) [\log (m r) -\log m)]^n + O(T^{\frac{1}{2}}\log^{n} T)\\
&= (-1)^{n} \sum_{m r \leq \frac{T}{2 \pi}} \Lambda (r) \log ^n r + O(T^{\frac{1}{2}}\log^{n} T).
\end{align*}
where the last three sums are over all integers $m$ and $r$ such that $mr \leq \frac{T}{2\pi}$. Our results follow from combining $J_1$ and $J_2$.
\end{proof}

\end{document}
